\chapter{Parallel CPU and GPU execution}
\label{sec:parallel-extension}

\begin{decisions}
\textbf{Status (26 October 2026).} The PAR-01 extension implements isolated
CPU/GPU workers, parallel endpoint repetitions, independent study replicas and
parallel population training/transfer. Development verification and the
independent closing audit are recorded in the master. The modified core is
\emph{not yet scientifically requalified}: PAR-QUAL-01 follows PAR-01, then the
mandatory GAP-01 final guard. Retained finite-local acceptance applies to its
original source, not automatically to this extension.
\end{decisions}

\section{What parallelism means}
The execution layer separates four decisions: concurrent objective calls,
method-valid batched proposals, independent complete studies and parallelism
inside a caller's training model. The first three have public routes here.
Giving one trial several GPUs makes those devices visible but does not create
DDP, model sharding or data loaders. Sequential adaptive methods preserve their
ask/tell dependency; their parallel route is independent replicas or endpoint
folds/seeds. Replicas remain separate experiments. Protected final evaluations
must never be used to choose the best replica retrospectively.

\section{Resource-aware workers}
\texttt{ResourceSpec(cpu\_cores=1, threads=1, gpus=0)} declares resources
\emph{per worker}. Counts are integers; threads cannot exceed reserved cores.
\texttt{ResourceExecutor(num\_workers, resources=..., cpu\_ids=...,
gpu\_ids=..., timeout=...)} validates capacity, partitions CPU affinity and
whole GPU devices, and keeps fresh-interpreter processes alive across calls.
Omitted CPU IDs use the coordinator's permitted affinity. GPU IDs must be
explicit physical indices or UUIDs; indices are canonicalized to UUIDs before
cooperative host leases are acquired. Use UUIDs for reproducible placement.

CUDA visibility and numerical thread environment variables are configured
before objective deserialization. Torch and threadpool limits are set in each
worker. CPU workers receive empty CUDA visibility. Objectives may declare an
\texttt{execution} parameter; its \texttt{ExecutionContext} contains worker,
task and attempt identities, CPU IDs, physical GPU UUIDs, thread count and the
logical device (\texttt{cpu} or \texttt{cuda:0}). Multi-device callers can use
logical indices within their assigned visible set. Returned trials include
worker identity and process metadata. Arbitrary mapped functions receive
\texttt{function(task, execution)} and return results in input order.

\begin{Verbatim}[fontsize=\small]
from hponas import ResourceExecutor, ResourceSpec, Config, Trial

def objective(config, execution):
    import torch
    x = torch.tensor(config['x'], dtype=torch.float64,
                     device=execution.device)
    return {'value': float((x - 0.375).square().cpu()),
            'cost': 1.0}

trials = [Trial(Config({'x': x}), str(i))
          for i, x in enumerate([0.1, 0.3, 0.5, 0.7])]
with ResourceExecutor(2, resources=ResourceSpec(), timeout=60) as ex:
    results = ex.execute_batch(trials, objective)
# GPU alternative: ResourceSpec(gpus=1), gpu_ids=[uuid_a, uuid_b].
# uuid_a and uuid_b must name available devices allocated to this run.
\end{Verbatim}

The local implementation requires Linux affinity, file locking and parent-death
signals. A persistent launcher thread owns worker creation so a replacement
worker survives the short-lived collection thread that requested it. A dead
coordinator terminates its workers; shutdown reclaims owned processes and
releases leases. Leases exclude other cooperating HPO-NAS executors, not
arbitrary applications. These are trusted Python workers, not a security
sandbox. Fractional GPU allocation and hard CPU/GPU memory partitions are
unsupported. Objective stdout/stderr are currently discarded by the worker
transport; application diagnostics should be returned as metadata or written
to caller-owned artifact paths.

\section{Study evaluation, budgets and recovery}
\texttt{Study(..., max\_concurrent=k)} controls outstanding configurations
subject to the searcher's capability contract. The separate
\texttt{evaluation\_workers=k} enables concurrent endpoint folds and seeds;
it defaults to one. Its capacity must fit both the executor and any
\texttt{StudyBudget.worker\_capacity}. Parallel declared streaming repetitions
are explicitly refused; ordinary streaming trial execution, reports,
coordinator stop decisions and cancellation remain supported.

\begin{Verbatim}[fontsize=\small]
# With a caller-defined StudyDeclaration and metered objective:
with ResourceExecutor(2, timeout=60) as ex:
    study = Study(searcher, ex, objective_fn,
        declaration=declaration, maximize=False,
        max_concurrent=1, evaluation_workers=2,
        resource_budget=StudyBudget(
            configurations=3, work_unit='updates',
            tuning_work=20, final_work=4, max_work_per_call=1,
            worker_capacity=2),
        checkpoint_dir='run-unique', study_name='declared')
    result = study.run()
\end{Verbatim}

An objective with this budget reports
\texttt{metadata['resource'] = \{'unit': 'updates', 'work': measured\}}.
The declaration owns disjoint tuning/final roles, folds, seeds and objective
identity. The coordinator durably reserves every admitted call before dispatch,
aggregates in fixed order and freezes selection before final evaluation.
Resource ledger version 2 records begin/complete/lost events in their actual
order: two calls admitted together remain legitimate if the first later
breaches a meter or elapsed limit. Recovery replays admission before completion,
validates raw measurements and rejects fabricated order or capacity.
Version 1 serial ledgers migrate through their original admission semantics.

Unknown interrupted work consumes the full declared per-call bound; measured
work is then a lower bound. No hidden retry is performed. An interrupted run
can continue only within its remaining admissible budget; lost time under a
time cap remains unknown. Endpoint time limits can overshoot by already
admitted calls. Worker timeouts are elapsed deadlines after worker admission,
not an end-to-end deadline including queue wait. Startup has its own deadline. A Study owns its executor and calls shutdown on completion or error; create a fresh executor for another Study. Independent replicas use the assigned no-op-shutdown wrapper inside the shared pool.

\section{Independent sequential-method replicas}
\begin{Verbatim}[fontsize=\small]
from hponas import run_replicas

def factory(seed, assigned_executor):
    # Make a fresh searcher and a different checkpoint path for each seed.
    return Study(make_searcher(seed), assigned_executor, objective_fn,
        declaration=declaration, budget=8, maximize=False,
        study_name=f'replica-{seed}',
        checkpoint_dir=f'replicas/{seed}').run()

with ResourceExecutor(2, timeout=120) as ex:
    independent_results = run_replicas(ex, factory, [7, 11])
\end{Verbatim}
\texttt{AssignedExecutor} executes inside an already placed worker and prevents
nested resource pools. The factory owns independent mutable state and paths.
The outer timeout bounds each complete replica task; size it for the bounded
study, not just one objective call. This interface does not synthesize a new
batch acquisition method or establish statistical superiority.

\section{Population and architecture transfer}
The classes \texttt{LocalPBT}, \texttt{LocalBGPBT} and \texttt{ArchitectureGeneration} share the resource pool when given \texttt{executor=ex}.
Without the optional executor they retain serial behavior. The supported
BG-PBT model remains the finite policy/value task with a shared
\texttt{CoordinateMLP} trunk; this change does not expand it to arbitrary PPO
simulators. The real finite-MDP objective, optimizer state and architecture
transfer execute on assigned devices.

The coordinator reserves an immutable population snapshot and ordered events,
then releases the population data lock during numerical work. Independent
members and transfer candidates run concurrently. Distillation rungs,
selection, temporal GP updates, exploitation and architecture-generation
barriers remain ordered. A separate lifetime coordinator lease prevents two
writers; it is distinct from the short data-journal lock. Results are saved
before publication, checked against target versions and replay-validated before
commit. Acknowledgement is idempotent. Dependent parent/child updates in the
same batch and stale writers are refused. Unknown failed batches retain their
full reservations. Serialized tensors and optimizer state remain CPU-portable.

The finite-MDP objective accepts both ordinary callables and Torch modules.
Module observations follow the first parameter or buffer; parameterless
callables receive the workload observations. Expectation tensors follow the
student output device, while teacher inputs follow the teacher's own placement.
Teacher outputs are explicitly detached before imitation loss construction.
This permits a GPU student with a CPU teacher and preserves the callable
known-answer references. The first PAR-QUAL-01 calibration exposed an
unconditional parameter lookup that broke these references; its failed
allocation is retained, and the compatibility repair is checked against the
original independent gradient and evidence-reconstruction tests.

\begin{Verbatim}[fontsize=\small]
from hponas.bg_pbt import LocalBGPBT
with ResourceExecutor(2, timeout=90) as ex:
    engine = LocalBGPBT('bg-run-unique', finite_objective,
        initial_members, architecture_pool,
        increment_bound=analytic_gain_bound,
        generations=2, max_rounds=1, patience=1,
        updates_per_round=1, distillation_rungs=(1, 2),
        max_updates=10, seed=7, fit_iterations=1, executor=ex)
    result = engine.run()
\end{Verbatim}
This small allocation assumes two members and four distinct architecture
candidates; the constructor derives and rejects insufficient update budgets.
The caller supplies valid objective/checkpoint/architecture objects as described
in the population guide. Reopening the same completed run returns its journal
without retraining.

\texttt{ParallelAdapterPopulation(..., executor=ex)} extends the model-neutral
adapter with ordered member/fold/seed batches and the same frozen final refit
schedule. GPU adapters must explicitly accept \texttt{execution} in
\texttt{initialize}, \texttt{restore}, \texttt{train}, \texttt{evaluate} and
\texttt{transfer}; they must place tensors on the assigned device and export
JSON/CPU-portable checkpoints. Missing GPU placement support is refused.
Partial uncertain adapter batches become an explicit terminal failure with
lower/upper work accounting; they are not silently retried.

\section{Algorithm routes and evidence}
\begin{longtable}{p{.26\textwidth}p{.34\textwidth}p{.31\textwidth}}
\toprule Family & Parallel route & Development evidence \\
\midrule\endhead
Random, Sobol, TPE, GP, CMA-ES & Independent same-seed replicas; existing pending/batch capabilities also govern Study & Analytic serial/replica traces; vector and scalar GPU fixtures \\
Robust GP, trust-region, cost GP, DEHB & Sequential updates inside independent replicas & Same-seed objective/configuration/fidelity traces \\
ASHA, PASHA, product/categorical PriorBand & Independent replicas preserve promotion and rung barriers & Same-seed scheduling traces and existing scheduler regressions \\
Finite guarded prior, $\pi$BO & Independent replicas preserve prior law/acquisition & Prior-route trace comparisons \\
Warm-start and ranked warm history & Separate replicas with frozen history & Seed queue, target identity and trace comparisons \\
RGPE & Parallel endpoint repetitions & Known-answer declared target; protected final selection \\
Vector random, qLogNEHVI, Chebyshev, NSGA-II & Existing batch proposals and concurrent trial evaluation & Actual two-worker analytic vector calls \\
PBT, PB2-Mix, BG-PBT & Independent members and transfer candidates within barriers & Exact CPU serial equality, recovery boundaries; real CUDA BG-PBT tolerance \\
Caller population adapter & Independent members and repetitions & Known-answer linear model, final refit, failure charging \\
\bottomrule
\end{longtable}
The executable roster lives under \path{tests/integration/}: \path{test_parallel_methods.py},
\path{test_parallel_population.py}, \path{test_parallel_study.py},
\path{test_resource_executor.py} and \path{test_parallel_hardware.py} under the
same test directory. Each receipt identifies its source and precise test count.
CPU evidence for a method is not evidence of that method's GPU execution.
ifBO remains the previously cancelled integration, not a newly supported family.
Legacy wrappers preserve their prior capability limits.

\section{Hardware observations and failure ledger}
The initial actual GPU check used two idle devices (physical indices 1 and 3,
recorded by UUID); unrelated jobs on devices 0 and 2 were untouched. Three
hardware tests passed in 21.24 seconds: isolated CUDA tensor evaluation against
an independent NumPy answer, CPU/CUDA BG-PBT with portable restore, and an
injected out-of-memory exception after a real CUDA allocation. The tiny tensor
fixture measured 0.0056793 seconds for one worker versus 0.00306712 for two:
speedup 1.85167 and efficiency 0.92584. These short measurements are descriptive,
not representative scaling results or an acceptance threshold. Subsequent
source repairs require new receipts; the initial measurements remain historical.

The current-source hardware successor passed all 12 tests in 83.79 seconds.
It additionally executes 16 scalar/scheduler/prior/warm families at two seeds
on real CUDA objectives, all four vector methods, RGPE endpoint repetitions,
PBT, PB2-Mix and a caller-owned CUDA linear adapter with portable checkpoints.
The same tensor timing now measures 0.00513631 versus 0.00311590 seconds
(speedup 1.64842; efficiency 0.82421), illustrating timing variability. Both
runs remain visible; neither is a representative benchmark. The complete
receipt is \path{20261025T190115Z-hardware-18860743/receipt.json} under the
PAR-01 evidence directory. Ranked warm-history selection was verified on CPU;
its GPU placement uses the same wrapper route but lacks a separate direct
GPU-specific selection test. Ray multi-GPU placement remains unverified.

The initial real local Ray test also passed with two logical CPU workers.
\texttt{RayExecutor} with a \texttt{ResourceSpec} and explicit \texttt{max\_retries=0} validates
cluster capacity and forwards execution context. Ray CPU reservations are
logical scheduler resources, not exclusive OS affinity. Ray GPU scheduling
uses its assigned visibility; local multi-GPU evidence does not qualify Ray
multi-GPU operation automatically.

Failures are preserved: initial recovery metadata tuple/list mismatch;
a vector-fixture schema mistake; sandbox restrictions on Ray sockets and GPU
inspection; and a development receipt invalidated by a mid-run test edit. Two expanded method-fixture runs also correctly refused a closed shared executor and overlapping CPU leases; ownership and disjoint test allocations were corrected before the passing 22-case rerun.
The boundary review repaired event-ordered budget replay, worker replacement
lifetime, pending population reservations, stale publication and process-ID
reuse in cleanup. Current-source reruns determine closure, never deletion of
these observations.

\section{Acceptance, limits and master execution}
\subsection{Critical-path qualification repair}
The third parallel qualification allocation completed all sixteen finite
campaigns and four capability gates, but its mutation run exposed missing
eligible coverage in the inherited \texttt{Study.run} grouped-dispatch branch.
Hardware-specific files were deliberately outside the single-CPU mutation
roster; that separation did not discharge the changed inherited critical path.
The run was stopped after concrete zero-return witnesses were retained. Its
disposition is \textbf{INCONCLUSIVE}, with 1,211 killed, 86 surviving, three
timed-out and 9,820 unchecked mutations; 21 critical survivors were unproved.
These partial counts are not a completed mutation score. The unchanged
zero-unproved-critical-survivor rule prevents acceptance. The stopped service
released its source freeze, and the publication waiter correctly refused to
publish a qualified manual.

The repair adds hardware-independent controller tests to the normal mutation
roster. On one CPU, real asynchronous endpoint calls must overlap across two
candidates, finish in reverse order and commit in declaration order. Analytic
scores, tuning-only selection, final-role separation, paid cancellation, lost
reservations and completed recovery are asserted. Separate tests execute the
actual Ray worker closures with a transport double and verify resource-spec
routing and execution-context propagation. Such tests establish controller
contracts; the independent real CPU/GPU tests still establish device execution.

A bounded diagnostic ran two clean controls and 28 exact retained mutations.
It detected 26 changes, including all original seven dispatch witnesses and
both Ray resource-dispatch changes. The two survivors only removed a redundant
cancellation argument. Parent cancellation is already normalized by
\texttt{Study.\_accept} before ledger publication, searcher observation and
scheduler reporting. The redundant argument is removed and scalar/vector
cancellation and recovery are checked. No new equivalence exclusion, weaker
threshold or reduced mutation population is introduced. The previous checker
and fixtures remain immutable; the successor binds the reviewed source while
preserving all 55 exact proof transformations.

This repair remains in PAR-QUAL-01. Current-source installation, regression,
isolation calibration and a prospectively reviewed qualification allocation
must pass before GAP-01/FINAL-AUDIT. Historical finite and hardware receipts
retain their original source identities. Neither this diagnostic nor a document
layout rehearsal qualifies the repaired source.

Exact CPU fixtures use exact equality. Prospective floating envelopes are
absolute/relative $10^{-9}/10^{-7}$ for float64 and $10^{-5}/10^{-4}$ for
float32. GPU timing synchronizes device work. Tests inspect placement, real
computation, ordering, accounting, restoration and failures; a missing GPU is
an unavailable result, never a pass. No global optimum, universal workload
speedup, arbitrary simulator compatibility or real-bank effectiveness is claimed.

PAR-01 follows CONTRACT, RESOURCES, EXECUTORS, STUDY, POPULATION, CONFORMANCE,
HARDWARE, DOCUMENTATION and AUDIT. At each boundary retain evidence, repair
observed defects, refresh the next bounded plan and continue. Development
rosters use four CPU cores at most, one numerical thread per worker, 300 seconds
per group and 8 GiB owned RSS; hardware uses at most two idle GPUs and 600
seconds. No sibling import, dataset or service is required. Separate-machine
testing remains waived. The final extension audit reconstructs source,
manifest, test, document and requirement coverage; PAR-QUAL-01 then addresses
current-source scientific qualification before GAP's terminal guard.

Authoritative plan: \path{docs/program/PARALLEL_EXECUTION_EXTENSION_2026_10_26.md}.
Exact operations: \path{docs/PROGRAM_OPERATIONS.md}.
Evidence and retained failures: \path{docs/program/evidence/2026-10-26/parallel-01/}.
